\chapter{引言}

\section{概述}

随着现代科学技术的发展和人们生活质量的提高，对于工程结构的性能提出
了越来越高的要求。例如：现代精密仪器、大型设备往往对于振动与位移有严格
的限制；生命线工程结构，要求在大震和大灾作用下依然保有必要的功能，以为
灾后救援与重建提供保障。20 世纪中叶以来，尽管社会发展水平有了巨大的提
高，然而由于灾害性作用而造成的损失却反而越来越大，这给结构工程学科带来
了一系列新的挑战性课题。正是在这样的背景下，基于性能的设计思想开始浮出
水面，并在近十年来引起了学者们强烈的兴趣\cite{chen2021}。


``自然界只有一个，自然现象遵循着不依赖于人类意志的客观规律。然而，
数理科学中却有着两套反映这些规律的体系：确定性描述和概率论描述。''虽
然概率论方法的发展引起了科学家和哲学家们关于自然本质的讨论，但是直到本
世纪五十年代以前，两套方法在各自独立的领域内都得到了长足的发展。六十年
代以来，由于本质非线性行为特别是混沌、分形等现象的发现和深入研究，随机
方法的重要性得到了日益深刻的认识。人们发现，在确定性非线性系统的长期
演化行为中会出现与随机行为不能加以区别的现象。而采用概率密度演化描述的
方法却能很好地描述其演化密度的长期行为\cite{chen2021,soave2020}。

\section{随机结构分析现状}

\subsection{线性随机结构分析}

经过三十余年的发展，线性随机结构在静力与动力分析方面的分析方法均已
趋于成熟。早期在物理学研究中使用的随机模拟方法于 20 世纪 70 年代初期引入
随机结构分析以来，已经成为检验各种随机结构分析方法的基本手段。基于随机
摄动展开的随机结构静力分析与动力分析也已于 20 世纪 80 年代基本完善。

......

......


本模板给出了典型样式的用例，供用户参考。样式用例包括公式、插图、子图、表格、参考文献等。样式用例见第~\ref{sec:style} 节。

\section{样式参考\label{sec:style}}
\subsection{公式}

公式居中，编号右对齐。如公式~\eqref{eq:example} 所示：
\begin{equation}
  E = mc^2 . \label{eq:example}
\end{equation}

多行公式可以采用 \texttt{align} 或者 \texttt{gather} 环境。
\begin{align}
  a &= b + c,  \label{eq:align1}\\
  d &= e + f + g,  \label{eq:align2}
\end{align}
或者
\begin{gather}
  a = b + c,  \label{eq:gather1}\\
  d = e + f + g.  \label{eq:gather2}
\end{gather}


插图建议先在外部绘制成 PDF 再引入，图题置于图下方。如图~\ref{fig:example} 所示。

\begin{figure}[htbp]
  \centering
  \includegraphics[width=0.72\linewidth,page=1]{demo-figures.pdf}
  \caption{单幅插图示例（位移时程）}
  \label{fig:example}
\end{figure}

含子图的插图用 \texttt{subcaption} 宏包的 \texttt{subfigure} 环境排版，各分图单独编号，
可以分别引用（图~\ref{fig:example-a}、图~\ref{fig:example-b}），如图~\ref{fig:subfigures} 所示。

\begin{figure}[htbp]
  \centering
  \begin{subfigure}[c]{0.48\linewidth}
    \centering
    \includegraphics[width=\linewidth,page=2]{demo-figures.pdf}
    \caption{位移时程}
    \label{fig:example-a}
  \end{subfigure}
  \hfill
  \begin{subfigure}[c]{0.48\linewidth}
    \centering
    \includegraphics[width=\linewidth,page=3]{demo-figures.pdf}
    \caption{幅值谱}
    \label{fig:example-b}
  \end{subfigure}
  \caption{含子图的插图示例}
  \label{fig:subfigures}
\end{figure}

表题通常置于表上方。表格统一用 \texttt{tabularray} 宏包排版，三线表写法如表~\ref{tab:example} 所示。
\begin{table}[htbp]
  \centering
  \caption{表格示例（三线表）}
  \label{tab:example}
  \begin{tblr}{
    colspec   = {c c c c},
    hline{1,Z} = {0.75bp},
    hline{2}   = {0.4bp},
    row{1}     = {font=\bfseries},
    rowsep     = 4pt,
  }
    模型   & 均值 & 标准差 & 变异系数 \\
    示例一 & 1.00 & 0.10   & 0.10     \\
    示例二 & 2.00 & 0.35   & 0.18     \\
    示例三 & 3.00 & 0.72   & 0.24     \\
  \end{tblr}
\end{table}

\section{参考文献}

以上示例给出了公式、插图、子图与表格的常用写法。图、表、公式均应在正文中先行引用，
再给出其内容；表格宽度一般不超过版心宽度，跨页的长表格可改用 \texttt{longtblr} 环境。
文献引用使用 \verb|\cite| 命令，参考文献列表由 \texttt{biblatex} 依据
\texttt{references.bib} 自动生成。
